% !TeX root = ../../main.tex
\subsection{英作文の問題について}
    英作文は北海道公立高校入試において、各大問の\fitblankbf[]{最後の問題}に用意されている。
    配点は英作文がメインの大問4で\fitblankbf[]{12}点。試験全体を通して大問4を含めて\fitblankbf[]{40}点分配点されている。
    そのうちのほとんどは、「\fitblankbf[]{主語と動詞を含む}1文で答えなさい」という形式で出題される。

    英作文は基本的に\fitblankbf[]{減点方式}で採点される。したがって、\uwave{単語や文法の間違いをなるべく避ける}
    ことが得点の鍵である。その具体的な方法はのちに示すこととする。

    大問4は3題あり、そのうち2題は\fitblankbf[]{空欄補充}、最後の1題が\fitblankbf[]{自由英作文}である。

\subsection{北海道公立高校入試の英作文における一般的な対策}
    リスニング同様、毎年、配点も問題構成も変わらない。仮に変わったとしても、全員が同じ問題を解くことになり、
    それを予想できた人間はいないはずであるから、次に示す対策をしない理由はない。

    \begin{enumerate}
        \item \fitblankbf[]{問題構成}をしっかりと理解する。
        \item \fitblankbf[]{24}語以上の英語で書かなければ、\fitblankbf[]{採点対象}にならない。
        \item 慣れてさえしまえば、長文読解に比べて、難易度的にも、時間的にも\fitblankbf[]{得点しやすい}。
    \end{enumerate}

\subsection{英作文における一般的な対策}
    まず、大前提として、英作文は\uwave{自分が表現できる英語で書く}ことが極めて重要である。
    具体的には、\uwave{書く内容が嘘}でも良いし、\uwave{幼稚な内容}だって構わない。
    なぜなら、採点の基準は\uwave{語数}と\uwave{単語、文法の正確さ}であり、内容が概ね本題に沿っていれば
    \uwave{内容の正しさ}は無関係であるからである。

    24語程度なら、多くの場合は2文〜3文程度で満たすことができる。

    特に、英作文をする上で必ず押さえておきたいポイントを次にまとめる。
\vspace{1em}

        \begin{techniquebox}{日々の学習で基礎語彙力、基礎文法力を身につけよう!}\label{technique:1}
            日々の学習で、\textbf{語彙の数}を増やし、\textbf{文法の精度}を高めることが大切である。
        \end{techniquebox}
\vspace{1em}

    単語が書けなければ英作文はできない。ここで勘違いして欲しくないのは、より高い語彙力や文法力を養うために、
    \uwave{より多くの単語、イディオム、慣用表現、文法を覚えなければならないということではない}。もちろん、
    それが達成できるに越したことはないが、英作文に必要な語彙はせいぜい\fitblankbf[]{100}語程度、
    文法は細かいものを含めても\fitblankbf[]{10}の文法程度であるかと思う。

    語彙力や文法力を高めることは、英作文の力を養うだけでなく、リスニング、長文読解など
    \fitblankbf[]{英語力全般}を高めるためにも大切である。
\newpage

        \begin{techniquebox}{戦える武器を持とう!}
            英作文において、\textbf{自分が得意とする語彙}、\textbf{自分が得意とする文法}
            を活用することが大切である。
        \end{techniquebox}
\vspace{1em}
    Technique\ref{technique:1}でも触れたが、英作文はにおいて、求められるものは良くも悪くも毎回同じようなものであるため、
    \uwave{自分が得意とする語彙や文法を活用する}ことが大切である。ここで、次に例をいくつか挙げる。

    \begin{enumerate}
        \item \fitblankbf[]{自分の意見と理由}を述べるとき
            \begin{itemize}
                \item 私は私の数学の先生は賢いと思います。
                \dialogueblank{I think that my math teacher is smart.}
                \item その理由は2つあります。
                \dialogueblank{There are [I have] two reasons for this.}
                \item 彼は、私をいつも幸せにしてくれます。
                \dialogueblank{He always makes me happy.}
                \item 私は彼に数学を教わっています。入試に合格するために。
                \dialogueblank{I am taught math by him to pass the entrance exam.}
            \end{itemize}
            \vspace{1em}
        \item \fitblankbf[]{副詞（句）と接続詞}を使うとき
            \par\vspace{0.5em}
            \noindent
            \begin{minipage}[t]{0.48\linewidth}
                \begin{itemize}
                    \item 一般的に
                    \dialogueblankcolumn{Generally,}
                    \item しかし
                    \dialogueblankcolumn{but , however}
                    \item さらに
                    \dialogueblankcolumn{even , moreover}
                    \item 一方で
                    \dialogueblankcolumn{on the other hand}
                \end{itemize}
            \end{minipage}\hfill
            \begin{minipage}[t]{0.48\linewidth}
                \begin{itemize}
                    \item 例えば
                    \dialogueblankcolumn{For example,}
                    \item 言い換えると
                    \dialogueblankcolumn{In other words,}
                    \item そのため
                    \dialogueblankcolumn{Therefore, as a result}
                    \item 最後に、
                    \dialogueblankcolumn{Finally,}
                \end{itemize}
            \end{minipage}
            \par\vspace{0.5em}
    \end{enumerate}
\vspace{1em}

    \begin{techniquebox}{日本語を工夫しよう!}
        \textbf{英語にしやすい日本語}で表現することが大切である。
    \end{techniquebox}
\vspace{1em}
英文を作る際、それが苦手な生徒ほど\uwave{日本語で考え込みすぎてしまう}ことが往々にしてある。
実は、しっかりとした日本語を用意すればするほど、それを英語に変換するのは難しくなる傾向にある。
例えば、以下の日本語を英作文するとき、どう考えるだろうか。
\newpage

\begin{itemize}
    \item 日本の技術は世界に貢献している。
    \dialogueblank{Japanese technology helps many people all over the world. \quad (contribute)}
    \item コミュニケーション能力を高めたい。
    \dialogueblank{I want to talk with many people well.}
    \item 英語がペラペラになりたい。（流暢に話したい）
    \dialogueblank{I want to speak English well. \quad (fluently)}
    \item 時間を無駄にすべきではない。
    \dialogueblank{We should use our time carefully.}
    \item 最後の試合に負けて、とても悲しかった。
    \dialogueblank{I was very sad because (when , that) we lost the last game.}
\end{itemize}
\vspace{1em}

このように、自分の表現したい日本語を\fitblankbf[]{英語にしやすい日本語}で表現することが大切である。
その際の具体的な方法として、\fitblankbf[]{言い換え}を用いることが大切である。
